\subsection{Consistencia residual de segundo orden}\label{sec:B4T05}
\begin{teorema}[Consistencia macroscópica]\label{thm:macro}
Fijemos $\delta_0,C,\eta,\gamma>0$. Sea $B$ un residual ordinario realizable
con $k\ge\delta_0n$ filas y columnas vacías. Supongamos que su punto de Bethe
$p$ es estacionario para $H_B$, interior, con masas RC uno y
\[
 p_e\le C/k,\qquad h_D\ge\eta,\qquad
 \lambda_2(\mathscr L_p)\ge\gamma.
\]
El grafo RC es conexo. El predictor usa el cofactor RC del determinante
canónico, el factor local y $\tfrac12\log(2\pi)$ con las convenciones de las Secciones~\ref{sec:B3T01} y~\ref{sec:B4T01}.
Entonces existen $k_0$ y $C_*$ dependientes sólo de estas cuatro constantes
tales que para $k\ge k_0$ cada hijo legal tiene un punto estacionario
interior único en la rama pequeña y
\[
 \boxed{\left|L(B+c)-L(B)-\log[p_c(1+\delta_c)]\right|
 \le C_*/k^2\quad\hbox{para todo hijo }c.}
\]
Aquí $\delta_c=(R\nabla\ell)_c/p_c=O(k^{-1})$.
No se excluyen hijos desfavorables ni se exige una cota inferior por casilla.
\end{teorema}
\begin{proof}
La existencia y el control del hijo se prueban en la Sección~\ref{sec:B4T02}, y la cancelación en las
Secciones~\ref{sec:B4T06}--\ref{sec:B4T09}, cuyo resultado final es la cota de $r_c$.
\end{proof}

\begin{corolario}\label{cor:rho}
En las condiciones del Teorema~\ref{thm:macro}, para cada fila vacía $r$,
$|\log\rho(B)|\le C_*/k^2$, donde $\rho(B)=\sum_{c\in r}e^{L(B+c)-L(B)}$.
\end{corolario}
\begin{proof}
La identidad diferencial de la Sección~\ref{sec:B3T01} da la suma uno en cada fila.
La positividad se sigue de $p_c>0$ y de $|\delta_c|=O(k^{-1})$ para $k$ grande (Sección~\ref{sec:B3T06}). Así,
\[
 w_c=p_c(1+\delta_c)>0,\qquad\sum_{c\in r}w_c=1,
 \quad \rho(B)=\sum_{c\in r}e^{L(B+c)-L(B)}=\mathbb E_w e^r.
\]
Por tanto $\min r\le\log\rho\le\max r$, lo que implica
\[
 \boxed{|\log\rho(B)|\le C_*/k^2.}
\]
\end{proof}
Si una trayectoria de padres permanece en una misma clase cuantitativa,
con las mismas constantes y $k\ge\delta_0n$, entonces
\[
 \sum_{k\ge\delta_0n}|\log\rho(B_k)|\le C_*\sum_{k\ge\delta_0n}k^{-2}
 =O_{\delta_0,C,\eta,\gamma}(n^{-1}).
\]
No se itera sin control la pérdida de parámetros de una herencia de un paso:
la pertenencia de toda la trayectoria es una hipótesis independiente.

\paragraph{Alcance.}
El teorema es determinista en la clase declarada. Que las trayectorias típicas
permanezcan en ella no se estudia aquí.
